%\needspace{4\baselineskip}
\color{uniblue}{\subsection{Блок команд управления}}
\color{black}


\begin{figure}[H]
\centering
\includegraphics[width=1\textwidth,height=1\textheight,keepaspectratio]{\fbpath/БКУ/img/BKU_logic.pdf}
\captionof{figure}{Перечень уставочных параметров команд управления}\label{pict:BKU_logic}
\end{figure}

\small
\setlength{\extrarowheight}{0.05cm} % высота ячеек в таблице
\begin{longtable}
{|>{\centering\arraybackslash}m{4.5cm}|>{\centering\arraybackslash}m{4.5cm}|>{\centering\arraybackslash}m{3cm}|>{\centering\arraybackslash}m{3.5cm}|}
\caption{Перечень уставочных параметров УВ\hfill\vspace{-0.5\baselineskip}}\label{ksn:tbl2} \\ 
\hhline{|-|-|-|-|}
\rowcolor{gray!30}
Наименование уставки & Диапазон & Шаг &  Значение по умолчанию \\ 
\hhline{|-|-|-|-|}
\endfirsthead

\caption*{\hspace{3pt}\emph{Продолжение таблицы \ref{ksn:tbl2}\hfill\vspace{-0.5\baselineskip}}} \\ % сделано по ГОСТ 2.105 п.6.8.7
\hhline{|-|-|-|-|}
\rowcolor{gray!30}
Наименование уставки & Диапазон & Шаг &  Значение по умолчанию \\  
\hhline{|-|-|-|-|}
\endhead
\hhline{|-|-|-|-|}

\endfoot
\endlastfoot

%%%%%%%%%%%%содержимое%%%%%%%%%%%%%%

\end{longtable}
\setlength{\extrarowheight}{0.0cm} % возврат параметра высота ячеек в таблице
\normalsize

\end{enumerate} 